% Source: https://artofproblemsolving.com/community/c3h2566942_fractanxfracpi4sin_xcos_x1
% Stable identity: aops-topic-2566942-q01
Using the angle-sum formulas,
\[\sin\left(x+\dfrac{\pi}{4}\right)=\dfrac{\sqrt2}{2}(\sin x+\cos x),\qquad \cos\left(x+\dfrac{\pi}{4}\right)=\dfrac{\sqrt2}{2}(\cos x-\sin x),\]
so $\tan\left(x+\dfrac{\pi}{4}\right)=\dfrac{\sin x+\cos x}{\cos x-\sin x}$, defined when $\cos x\ne\sin x$. The equation also needs $\sin x+\cos x\ne0$. Under these conditions it becomes
\[\dfrac{1}{\cos x-\sin x}=1\iff\cos x-\sin x=1.\]
On the unit circle $(X,Y)=(\cos x,\sin x)$ this is the line $X-Y=1$, which meets the circle at $(1,0)$ and $(0,-1)$ (substituting $X=1+Y$ into $X^2+Y^2=1$ gives $Y(Y+1)=0$). So $x=2k\pi$ or $x=\dfrac{3\pi}{2}+2k\pi$. For these, $\sin x+\cos x=\pm1\ne0$ and $\cos x-\sin x=1\ne0$, and indeed $\dfrac{\tan\dfrac{\pi}{4}}{1}=1$ and $\dfrac{\tan\left(-\dfrac{\pi}{4}\right)}{-1}=1$.

The solutions are
\[\boxed{x=2k\pi\quad\text{or}\quad x=\dfrac{3\pi}{2}+2k\pi,\qquad k\in\mathbb Z.}\]

\begin{asy}
usepackage("amsmath");
size(4.5cm);
defaultpen(fontsize(10pt));
draw(circle((0,0),1));draw((-1.4,0)--(1.6,0),Arrow(TeXHead));draw((0,-1.4)--(0,1.4),Arrow(TeXHead));
draw((-0.25,-1.25)--(1.3,0.3),linewidth(0.8));
dot((1,0));dot((0,-1));
label("$x=0$",(1,0),NW,fontsize(9pt));label("$x=\dfrac{3\pi}{2}$",(0,-1),SE,fontsize(9pt));
label("$X-Y=1$",(1.3,0.3),NE,fontsize(9pt));
label("$X$",(1.6,0),E);label("$Y$",(0,1.4),W);
\end{asy}
